\begin{table}[t]
\centering
\small
\begin{tabular}{lcccc}
\toprule
Task   & \molm & \smi & Random \\
\midrule
atom type        & 1.00 & 1.00 & 0.24 \\
is aromatic       & 1.00 & 1.00 & 0.48 \\
is in ring        & \textbf{0.85} & \textbf{1.00} & 0.49 \\
degree           & \textbf{0.62} & \textbf{0.99} & 0.24 \\
hybridization    & \textbf{0.86} & \textbf{0.99} & 0.30 \\
chiral tag      & 0.52 & 0.94 & 0.31 \\
formal charge    & 0.83 & 1.00 & 0.33 \\
total valence     & 0.96 & 1.00 & 0.23 \\
\bottomrule
\end{tabular}
\caption{\textbf{Best-layer linear-probe macro-F1 per task.} The cross-model gap is largest on contextual structural properties (\textbf{bold}: ring membership, degree, hybridization), which cannot be derived from the SMILES token alone. Token-determined tasks (atom type, aromaticity) saturate for both models; chirality and formal charge are tokenizer-determined for \smi (bracket-atom vocabulary) but require contextual reconstruction by \molm. ``Majority F1'' is the macro-F1 of a constant predictor that always returns the majority class; ``Random'' is a linear probe trained on Gaussian features of the same shape.}
\label{tab:probing-headline}
\end{table}